\chapter{P5: Distribusi Event}
\label{chap:p5}

\section{Rancangan Solusi}

Aggregator memublikasikan perubahan hazard melalui exchange RabbitMQ bertipe
fanout. Setiap consumer memiliki queue sendiri, sehingga satu publication
disalin ke Event Live Feed, Webhook Provider, dan Danger Notification.
Ketiga layanan dapat memproses publication yang sama untuk tujuan berbeda.
Menempatkan ketiganya pada satu queue akan menghasilkan pembagian pekerjaan
antar-consumer, bukan penyalinan event ke seluruh layanan.

\begin{table}[htbp]
    \centering
    \small
    \caption{Topologi distribusi event}
    \begin{tabularx}{\textwidth}{@{}P{0.22\textwidth}X P{0.25\textwidth}@{}}
        \toprule
        \textbf{Komponen} & \textbf{Nama contoh/default} & \textbf{Peran} \\
        \midrule
        Exchange & \code{bnpb.hazards.v1} & Fanout durable. \\
        Queue feed & \code{bnpb.event-live-feed.v1} & Event streaming dan replay. \\
        Queue webhook & \code{bnpb.webhook-provider.v1} & Callback subscription. \\
        Queue notification & \code{bnpb.danger-notification.v1} & Push SIAGA/AWAS. \\
        DLQ per consumer & Nama layanan dengan akhiran \code{.dlq.v1} & Menampung pesan yang ditolak tanpa requeue. \\
        \bottomrule
    \end{tabularx}
\end{table}

Nama exchange dan queue berasal dari konfigurasi. Producer hanya
mendeklarasikan exchange; tiap consumer mendeklarasikan queue durable,
dead-letter queue, dan binding miliknya. Routing dead-letter menggunakan
default exchange dengan routing key nama DLQ. Ketiga consumer menggunakan
prefetch satu pesan dan acknowledgement manual. Pengaturan ini membatasi
pesan yang belum diakui, tetapi juga membuat pekerjaan lambat dalam satu
consumer menahan pekerjaan berikutnya pada consumer tersebut.

Publication memakai envelope Protobuf \code{HazardEventPublished}, berisi
\code{event_id}, \code{schema_version}, sumber, publisher, correlation ID,
waktu publication, dan payload hazard kanonis. Implementasi saat ini
menggunakan versi schema 1 dan publisher \code{aggregator}.
\code{event_id} berupa UUID publication dan berbeda dari \code{hazard_id}
yang stabil untuk satu referensi sumber. Pembentukan envelope dapat
ditelusuri pada \code{services/aggregator/internal/ingest/mapper.go}.

Alur penyimpanan dan publication menggunakan pola transactional outbox:
\begin{enumerate}
    \item Poller memetakan data sumber ke hazard kanonis, lalu membuka
    transaksi PostgreSQL.
    \item Upsert hazard dan penyisipan envelope ke \code{outbox_events}
    menggunakan transaksi yang sama, kemudian transaksi di-commit.
    \item Publisher mengambil paling banyak 20 publication pending
    setiap 200 ms dan mengirim pesan persisten ke exchange.
    \item Publisher menunggu publisher confirm dari broker sebelum
    mengisi \code{published_at}. Kegagalan publish menaikkan retry record.
    \item Pada batas retry, contoh konfigurasi lima percobaan, record
    ditandai gagal dan disalin ke tabel \code{outbox_dead_letters}.
\end{enumerate}

Implementasi publisher terdapat pada
\code{services/aggregator/internal/broker/publisher.go} dan
\code{services/aggregator/internal/store/outbox_store.go}.
Tabel dead-letter outbox merupakan state PostgreSQL producer, berbeda
dari DLQ RabbitMQ masing-masing consumer. Kegagalan koneksi awal broker
diulang setiap dua detik; retry koneksi tersebut bukan penghitung percobaan
publish setiap record.

Setiap consumer menentukan titik acknowledgement sesuai pekerjaannya:
\begin{itemize}
    \item \textbf{Event Live Feed} memvalidasi envelope, menambahkan event
    ke replay buffer, lalu ack. ID yang masih berada dalam replay window
    dideduplikasi. Envelope rusak atau tidak didukung langsung di-nack
    tanpa requeue. Koneksi diulang dengan backoff 1--30 detik.
    \item \textbf{Webhook Provider} mencari subscription yang sesuai filter,
    memeriksa dedup, membentuk payload sesuai view, dan mengirim HTTP POST
    dengan tanda tangan HMAC-SHA256 pada \code{X-BNPB-Signature}.
    Hasil dicatat ke SQLite dan delivery sukses ditandai untuk dedup.
    Default dispatcher melakukan lima percobaan dengan timeout 10 detik
    per request serta jeda 1, 2, 4, dan 8 detik.
    \item \textbf{Danger Notification} memproses envelope, memilih severity
    SIAGA/AWAS, kemudian mengirim melalui topic FCM. Ack dilakukan setelah
    pemrosesan berhasil. Parse error Protobuf langsung ditolak; error
    lain memperoleh satu redelivery sebelum ditolak menuju DLQ.
\end{itemize}

Kode consumer berada pada
\code{services/event-live-feed/src/adapters/inbound/rabbitmq.rs},
\code{services/webhook-provider/internal/consumer/consumer.go}, dan
\code{services/danger-notification/BNPB.Notification.Server/Adapters/Inbound/RabbitMqConsumer.cs}.
Header trace dibawa melalui AMQP untuk menghubungkan span publish dan
consume. Identitas publication membantu menelusuri satu event di beberapa
consumer, sementara correlation ID menghubungkannya dengan putaran ingestion.

\section{Justifikasi dan Alternatif}

Fanout dipilih karena semua consumer menerima event hazard yang sama lalu
melakukan filter atau proyeksi sesuai fungsi layanan. Producer tidak perlu
mengetahui alamat callback maupun perangkat pengguna. Queue terpisah
memungkinkan consumer yang berhenti menimbun backlog tanpa memaksa consumer
lain berhenti memproses queue masing-masing, selama broker dan resource
bersama masih tersedia.

Alternatif panggilan HTTP langsung dari Aggregator ke setiap consumer
mengurangi kebutuhan broker, tetapi menjadikan publisher bertanggung jawab
atas timeout, retry, dan availability setiap tujuan. Alternatif exchange
topic dapat mengurangi lalu lintas dengan routing menurut sumber atau jenis
hazard; konsekuensinya ialah penambahan aturan routing dan kesepakatan key.
Fanout mempertahankan topologi yang sederhana untuk tiga consumer saat ini.

Alternatif log terdistribusi berbasis offset memberi retensi dan replay
yang lebih panjang, tetapi memerlukan pengelolaan partition, offset, dan
consumer group. Kebutuhan milestone dapat ditangani RabbitMQ dan outbox
PostgreSQL. Keputusan ini tidak menyamakan queue dengan arsip event permanen:
event yang telah di-ack tidak lagi menjadi backlog broker, dan replay feed
memiliki retensi sendiri.

\section{Batasan Solusi}

Pada transport broker, persistent message, publisher confirm, manual ack,
dan redelivery mendukung pola \emph{at-least-once} selama queue sudah terikat,
state broker masih tersedia, dan pesan belum dihentikan oleh kebijakan retry
atau DLQ. Confirm producer dan ack consumer mengonfirmasi tahap berbeda~\cite{rabbit}.
Jaminan tersebut tidak sama dengan keberhasilan efek bisnis untuk setiap
subscriber. Semantik sistem belum exactly-once. Publisher confirm dan penandaan
\code{published_at} merupakan dua operasi berbeda, sehingga crash di
antaranya dapat menyebabkan publication dikirim ulang. Demikian pula,
keberhasilan HTTP/FCM dan acknowledgement broker tidak berada dalam
transaksi yang sama. Mekanisme dedup harus dinilai per consumer:
\begin{enumerate}
    \item Feed hanya mendeduplikasi sepanjang replay window dalam memori.
    Ack setelah append tidak berarti data tersimpan di disk; restart
    menghapus replay dan himpunan dedup.
    \item Webhook menyimpan dedup persisten per pasangan publication dan
    subscription, tetapi crash setelah penerima menerima POST dan sebelum
    penandaan dedup masih dapat menghasilkan pengiriman ulang. Penerima
    juga perlu menangani duplikasi.
    \item Notification belum menyimpan dedup publication atau riwayat
    pengiriman. Redelivery setelah FCM menerima pesan masih dapat
    menghasilkan push berulang.
\end{enumerate}

Retry webhook memiliki keterbatasan khusus: setelah seluruh percobaan HTTP
gagal, consumer mengatur kegagalan pemrosesan dan melakukan nack dengan
requeue. Belum ada batas total redelivery pada jalur tersebut. Dengan
demikian, lima kegagalan dispatcher tidak otomatis memindahkan event ke
DLQ. Komentar kode yang menyebut dispatch failure tidak menghalangi ack
tidak sesuai dengan alur kondisi aktual; uraian laporan mengikuti alur
yang dieksekusi.

Pola outbox juga belum menutup semua jalur kegagalan aplikasi. Poller
mengabaikan error pembentukan envelope dan hasil pemanggilan
\code{EnqueueOutbox}. Transaksi SQL yang sama mendukung konsistensi ketika
operasinya berhasil, tetapi belum cukup untuk menjanjikan bahwa setiap
hazard pasti mempunyai publication pada seluruh jalur error. Pembaruan
watermark dan publikasi berulang perlu diuji secara terpisah.

Publication memakai \code{mandatory=false}; apabila queue consumer belum
pernah dibuat dan diikat ketika pesan dipublikasikan, publisher confirm
tidak membuktikan consumer tersebut memperoleh salinannya. Compose belum
memasang volume broker secara eksplisit. Karena itu, label durable pada
queue dan persistent pada pesan tidak cukup untuk menjanjikan pemulihan
setelah container broker dibuat ulang. Pemeriksaan schema juga tidak
seragam: feed memvalidasi publisher, ID, sumber, dan payload, sedangkan
notification terutama bergantung pada parsing serta akses field yang dipakai.

\section{Proof-of-Concept dan Keterlacakan}

Pada snapshot 9 Oktober 2026 tersedia publisher, tiga consumer, kontrak
event, serta skenario pengujian. Belum terdapat hasil pengujian fanout,
redelivery, dan pemulihan menyeluruh terbaru yang dilampirkan. Bukti
container mock historis tidak digunakan sebagai pengganti bukti
RabbitMQ dan ketiga consumer.

\begin{table}[htbp]
    \centering
    \small
    \caption{Artefak pengujian dan batas cakupannya}
    \begin{tabularx}{\textwidth}{@{}P{0.28\textwidth}X@{}}
        \toprule
        \textbf{Artefak} & \textbf{Cakupan yang dapat ditelusuri} \\
        \midrule
        \code{tests/test_e2e_feeds.py} & Penerimaan event SSE, filter seismic/volcanic, dan izin full view. \\
        \code{tests/test_e2e_webhooks.py} & Lifecycle subscription, tanda tangan ping/event, serta pencatatan delivery. \\
        \code{tests/test_e2e_notifications.py} & Pembacaan notification dan lifecycle token; tidak dengan sendirinya membuktikan dedup push. \\
        Tes unit adapter consumer & Feed menguji validasi envelope; notification menguji keputusan requeue. Keberadaan tes bukan hasil eksekusi terbaru. \\
        \bottomrule
    \end{tabularx}
\end{table}

Untuk pembuktian yang dapat diulang, ketiga consumer harus dinyalakan lebih
dahulu hingga queue dan binding tersedia. Catat satu \code{event_id} yang
dipublikasikan Aggregator, lalu cocokkan penerimaan pada feed, log delivery
webhook, dan pemrosesan notification. Gunakan hazard SIAGA/AWAS apabila
observasi mencakup pengiriman FCM. Simpan hasil publisher confirm, ack,
queue depth, dan trace terkait tanpa menyalin token atau secret.

Skenario berikutnya menghentikan satu consumer setelah binding terbentuk,
memublikasikan event, dan memeriksa bahwa consumer lain tetap memproses
salinannya. Setelah consumer dinyalakan kembali, periksa backlog dan
redelivery pada queue miliknya. Uji pesan tidak valid secara terpisah untuk
memastikan arah DLQ. Uji replay \code{event_id} yang sama harus dibedakan
dari upsert baru dengan \code{event_id} berbeda, karena keduanya mempunyai
makna deduplikasi yang berbeda.

Eksekusi E2E mengikuti \code{docs/e2e-testing.md}, dengan stack aktif dan
\code{BNPB_E2E=1}. Hasil yang perlu dilengkapi mencakup penerimaan di tiga
consumer, jumlah pengiriman ulang, perilaku DLQ, dan kehilangan replay saat
restart feed. Sampai hasil tersebut direkam, status pembuktian P5 adalah
mekanisme dan skenario tersedia, sedangkan keberhasilan integrasi serta
karakteristik pemulihan terbaru belum diukur.

Untuk menunjukkan penambahan consumer tanpa perubahan producer, buat queue
pengamat baru dengan nama unik dan bind ke exchange fanout yang sama melalui
RabbitMQ Management atau consumer uji. Setelah binding aktif, tunggu publication
baru dan cocokkan \code{event_id} dengan minimal dua queue layanan yang sudah
ada. Rekam konfigurasi binding, payload yang telah disanitasi, dan identitas
container Aggregator yang tetap berjalan. Consumer baru menerima publication
sejak binding dibuat, bukan otomatis seluruh sejarah sebelum berlangganan.
Skenario ini belum dieksekusi sebagai bagian bukti laporan.
